% steps/step02.tex
\section[گام دوم: تحلیل نیازمندی‌ها]{\faTasks\ \ گام دوم: تحلیل نیازمندی‌ها}
\begin{stepbox}

\field{کاربران و ذینفعان سیستم چه کسانی هستند؟}{
    کاربران اصلی این سیستم، رادیولوژیست‌ها و پزشکان متخصص آنکولوژی هستند که می‌توانند از آن به‌عنوان یک دستیار هوشمند برای افزایش سرعت و دقت تشخیص بهره ببرند. علاوه بر این، بیماران نیز به‌صورت غیرمستقیم از تشخیص زودهنگام تومورها، کاهش خطای انسانی و جلوگیری از تشخیص‌های مثبت کاذب منتفع خواهند شد.
}

\field{دسترسی به داده‌ها و منابع چگونه است؟}{
    مهم‌ترین چالش در این بخش، حفظ حریم خصوصی بیماران است که دسترسی آزاد به تصاویر پزشکی را محدود می‌کند. از سوی دیگر، داده‌های موجود معمولاً در قالب‌های مختلف ذخیره شده‌اند و پیش از ورود به مدل، نیازمند استانداردسازی، یکپارچه‌سازی فرمت و آماده‌سازی ساختاری هستند. بنابراین کیفیت و سازمان‌دهی داده‌ها نقش تعیین‌کننده‌ای در موفقیت پروژه دارد.
}

\field{محدودیت‌های فنی، مالی، زیرساختی و زمانی چیست؟}{
    از منظر فنی، تصاویر پزشکی معمولاً به‌صورت سه‌بعدی و در فرمت‌هایی مانند NIfTI ذخیره می‌شوند و پردازش آن‌ها به منابع سخت‌افزاری قدرتمند، به‌ویژه RAM و GPU بالا، نیاز دارد. در نتیجه، استفاده از روش‌هایی مانند پردازش قطعه‌به‌قطعه (Patch-based) برای مدیریت بار محاسباتی ضروری است. همچنین، داده‌های CT Scan بر حسب واحد HU مقداردهی می‌شوند و پیاده‌سازی صحیح تکنیک Windowing برای استخراج بازه مناسب بافت کبد، یک چالش فنی مهم به‌شمار می‌رود. از نظر زمانی نیز، به دلیل حساسیت بالای کاربردهای پزشکی و نیاز به اعتبارسنجی‌های تخصصی و قانونی، چرخه توسعه چنین سامانه‌ای بسیار طولانی‌تر از پروژه‌های معمول پردازش تصویر است و ممکن است از مرحله ایده تا استقرار بالینی حدود ۱ تا ۳ سال زمان ببرد.
}

\end{stepbox}
